\begingroup
\fontsize{8.6}{10.5}\selectfont
\begin{thebibliography}{126}
\bibitem{ref001}
Hosack, T., Damry, D., and Biswas, S. (2023). Drug-induced liver injury: a comprehensive review. Ther. Adv. Gastroenterol. \textit{16}, 17562848231163410. \url{https://doi.org/10.1177/17562848231163410}.

\bibitem{ref002}
Tiwari, V., Shandily, S., Albert, J., Mishra, V., Dikkatwar, M., Singh, R., Sah, S.K., and Chand, S. (2025). Insights into medication-induced liver injury: Understanding and management strategies. Toxicol. Rep. \textit{14}, 101976. \url{https://doi.org/10.1016/j.toxrep.2025.101976}.

\bibitem{ref003}
Chidiac, A.S., Buckley, N.A., Noghrehchi, F., and Cairns, R. (2023). Paracetamol (acetaminophen) overdose and hepatotoxicity: mechanism, treatment, prevention measures, and estimates of burden of disease. Expert Opin. Drug Metab. Toxicol. \textit{19}, 297--317. \url{https://doi.org/10.1080/17425255.2023.2223959}.

\bibitem{ref004}
Yoon, E., Babar, A., Choudhary, M., Kutner, M., and Pyrsopoulos, N. (2016). Acetaminophen-Induced Hepatotoxicity: a Comprehensive Update. J. Clin. Transl. Hepatol. \textit{4}, 131--142. \url{https://doi.org/10.14218/JCTH.2015.00052}.

\bibitem{ref005}
Ramachandran, A., and Jaeschke, H. (2019). Acetaminophen Hepatotoxicity. Semin. Liver Dis. \textit{39}, 221--234. \url{https://doi.org/10.1055/s-0039-1679919}.

\bibitem{ref006}
Bhushan, B., and Apte, U. (2019). Liver Regeneration after Acetaminophen Hepatotoxicity: Mechanisms and Therapeutic Opportunities. Am. J. Pathol. \textit{189}, 719--729. \url{https://doi.org/10.1016/j.ajpath.2018.12.006}.

\bibitem{ref007}
Aoyagi, T., Goya, T., Imoto, K., Azuma, Y., Hioki, T., Kohjima, M., Tanaka, M., Oda, Y., and Ogawa, Y. (2025). Two types of regenerative cell populations appear in acute liver injury. Stem Cell Reports \textit{20}, 102503. \url{https://doi.org/10.1016/j.stemcr.2025.102503}.

\bibitem{ref008}
Kotulkar, M., Paine-Cabrera, D., Abernathy, S., Robarts, D.R., Parkes, W.S., Lin-Rahardja, K., Numata, S., Lebofsky, M., Jaeschke, H., and Apte, U. (2023). Role of HNF4alpha-cMyc interaction in liver regeneration and recovery after acetaminophen-induced acute liver injury. Hepatology \textit{78}, 1106--1117. \url{https://doi.org/10.1097/HEP.0000000000000367}.

\bibitem{ref009}
Umbaugh, D.S., Ramachandran, A., and Jaeschke, H. (2021). Spatial Reconstruction of the Early Hepatic Transcriptomic Landscape After an Acetaminophen Overdose Using Single-Cell RNA-Sequencing. Toxicol. Sci. \textit{182}, 327--345. \url{https://doi.org/10.1093/toxsci/kfab052}.

\bibitem{ref010}
Walesky, C.M., Kolb, K.E., Winston, C.L., Henderson, J., Kruft, B., Fleming, I., Ko, S., Monga, S.P., Mueller, F., Apte, U., et al. (2020). Functional compensation precedes recovery of tissue mass following acute liver injury. Nat. Commun. \textit{11}, 5785. \url{https://doi.org/10.1038/s41467-020-19558-3}.

\bibitem{ref011}
Ben-Moshe, S., Veg, T., Manco, R., Dan, S., Papinutti, D., Lifshitz, A., Kolodziejczyk, A.A., Bahar Halpern, K., Elinav, E., and Itzkovitz, S. (2022). The spatiotemporal program of zonal liver regeneration following acute injury. Cell Stem Cell \textit{29}, 973--989.e10. \url{https://doi.org/10.1016/j.stem.2022.04.008}.

\bibitem{ref012}
Hutton, A., and Meyer, J.G. (2025). Trajectory Inference for Single Cell Omics. Preprint at arXiv. \url{https://doi.org/10.48550/arXiv.2502.09354}.

\bibitem{ref013}
Deconinck, L., Cannoodt, R., Saelens, W., Deplancke, B., and Saeys, Y. (2021). Recent advances in trajectory inference from single-cell omics data. Curr. Opin. Syst. Biol. \textit{27}, 100344. \url{https://doi.org/10.1016/j.coisb.2021.05.005}.

\bibitem{ref014}
La Manno, G., Soldatov, R., Zeisel, A., Braun, E., Hochgerner, H., Petukhov, V., Lidschreiber, K., Kastriti, M.E., Lönnerberg, P., Furlan, A., et al. (2018). RNA velocity of single cells. Nature \textit{560}, 494--498. \url{https://doi.org/10.1038/s41586-018-0414-6}.

\bibitem{ref015}
Bergen, V., Lange, M., Peidli, S., Wolf, F.A., and Theis, F.J. (2020). Generalizing RNA velocity to transient cell states through dynamical modeling. Nat. Biotechnol. \textit{38}, 1408--1414. \url{https://doi.org/10.1038/s41587-020-0591-3}.

\bibitem{ref016}
Forrow, A., and Schiebinger, G. (2021). LineageOT is a unified framework for lineage tracing and trajectory inference. Nat. Commun. \textit{12}, 4940. \url{https://doi.org/10.1038/s41467-021-25133-1}.

\bibitem{ref017}
Sashittal, P., Zhang, R.Y., Law, B.K., Schmidt, H., Strzalkowski, A., Bolondi, A., Chan, M.M., and Raphael, B.J. (2026). Inferring cell differentiation maps from lineage tracing data. Nat. Methods \textit{23}, 532--541. \url{https://doi.org/10.1038/s41592-025-02903-z}.

\bibitem{ref018}
Bhushan, B., Gunewardena, S., Edwards, G., and Apte, U. (2020). Comparison of liver regeneration after partial hepatectomy and acetaminophen-induced acute liver failure: A global picture based on transcriptome analysis. Food Chem. Toxicol. \textit{139}, 111186. \url{https://doi.org/10.1016/j.fct.2020.111186}.

\bibitem{ref019}
Matchett, K.P., Wilson-Kanamori, J.R., Portman, J.R., Kapourani, C.A., Fercoq, F., May, S., Zajdel, E., Beltran, M., Sutherland, E.F., Mackey, J.B.G., et al. (2024). Multimodal decoding of human liver regeneration. Nature \textit{630}, 158--165. \url{https://doi.org/10.1038/s41586-024-07376-2}.

\bibitem{ref020}
Yu, L., Yan, J., Zhan, Y., Li, A., Zhu, L., Qian, J., Zhou, F., Lu, X., and Fan, X. (2023). Single-cell RNA sequencing reveals the dynamics of hepatic non-parenchymal cells in autoprotection against acetaminophen-induced hepatotoxicity. J. Pharm. Anal. \textit{13}, 926--941. \url{https://doi.org/10.1016/j.jpha.2023.05.004}.

\bibitem{ref021}
Liu, H., Zhang, Y., and Ning, S. (2025). ScRNA-seq combined with ATAC-seq analysis to explore the metabolic balance mechanism of CCl4-induced liver inflammatory injury. Front. Immunol. \textit{16}, 1600685. \url{https://doi.org/10.3389/fimmu.2025.1600685}.

\bibitem{ref022}
Bwanya, B., Verheijen, M.C.T., Hauser, D., de Kok, T.M., Jennen, D.G.J., van den Beucken, T., and Caiment, F. (2026). Single-cell transcriptomics of acetaminophen-induced responses in human 2D and 3D liver microtissues. Arch. Toxicol. \textit{100}, 1411--1424. \url{https://doi.org/10.1007/s00204-025-04296-6}.

\bibitem{ref023}
Ravindra, N.G., Sehanobish, A., Pappalardo, J.L., Hafler, D.A., and van Dijk, D. (2020). Disease State Prediction From Single-Cell Data Using Graph Attention Networks. In Proceedings of the ACM Conference on Health, Inference, and Learning (ACM), pp. 121--130. \url{https://doi.org/10.1145/3368555.3384449}.

\bibitem{ref024}
Wang, J., Ma, A., Chang, Y., Gong, J., Jiang, Y., Qi, R., Wang, C., Fu, H., Ma, Q., and Xu, D. (2021). scGNN is a novel graph neural network framework for single-cell RNA-Seq analyses. Nat. Commun. \textit{12}, 1882. \url{https://doi.org/10.1038/s41467-021-22197-x}.

\bibitem{ref025}
Zhang, M.J., Hou, K., Dey, K.K., Sakaue, S., Jagadeesh, K.A., Weinand, K., Taychameekiatchai, A., Rao, P., Pisco, A.O., Zou, J., et al. (2022). Polygenic enrichment distinguishes disease associations of individual cells in single-cell RNA-seq data. Nat. Genet. \textit{54}, 1572--1580. \url{https://doi.org/10.1038/s41588-022-01167-z}.

\bibitem{ref026}
Zhang, S., Shu, H., Zhou, J., Rubin-Sigler, J., Yang, X., Liu, Y., Cooper-Knock, J., Monte, E., Zhu, C., Tu, S., et al. (2026). Single-cell polygenic risk scores dissect cellular and molecular heterogeneity of complex human diseases. Nat. Biotechnol. \textit{44}, 845--861. \url{https://doi.org/10.1038/s41587-025-02725-6}.

\bibitem{ref027}
Liu, C., Yuan, B., Shen, B., Li, J., Zhu, R., Yang, P., Wu, B., Xuan, Y., Yang, S., Yang, N., et al. (2026). DPCGS: a computational framework for linking GWAS to single-cell transcriptomics in complex traits and diseases. Preprint at bioRxiv (version 2). \url{https://doi.org/10.64898/2026.07.14.738331}.

\bibitem{ref028}
Nandhini, P., and Vimalnath, S. (2025). Optimized multi scale graph neural network with attention mechanism for cooperative spectrum sensing in cognitive radio networks. Sci. Rep. \textit{15}, 41130. \url{https://doi.org/10.1038/s41598-025-24947-z}.

\bibitem{ref029}
Nabian, M.A., Liu, C., Ranade, R., and Choudhry, S. (2024). X-MeshGraphNet: Scalable Multi-Scale Graph Neural Networks for Physics Simulation. Preprint at arXiv. \url{https://doi.org/10.48550/arXiv.2411.17164}.

\bibitem{ref030}
Barwey, S., Pal, P., Patel, S., Balin, R., Lusch, B., Vishwanath, V., Maulik, R., and Balakrishnan, R. (2025). Mesh-based super-resolution of fluid flows with multiscale graph neural networks. Comput. Methods Appl. Mech. Eng. \textit{443}, 118072. \url{https://doi.org/10.1016/j.cma.2025.118072}.

\bibitem{ref031}
Gal, E., Eliasof, M., Schönlieb, C.-B., Kyrchei, I.I., Haber, E., and Treister, E. (2025). Towards Efficient Training of Graph Neural Networks: A Multiscale Approach. Trans. Mach. Learn. Res. \url{https://openreview.net/forum?id=2eZ8xkL2ZB}.

\bibitem{ref032}
Fan, W., Fei, J., Guo, D., Yi, K., Song, X., Xiang, H., Ye, H., and Li, M. (2025). MedGNN: Towards Multi-resolution Spatiotemporal Graph Learning for Medical Time Series Classification. In Proceedings of the ACM Web Conference 2025 (ACM), pp. 5054--5064. \url{https://doi.org/10.1145/3696410.3714514}.

\bibitem{ref033}
Jaeschke, H., Williams, C.D., Ramachandran, A., and Bajt, M.L. (2012). Acetaminophen hepatotoxicity and repair: the role of sterile inflammation and innate immunity. Liver Int. \textit{32}, 8--20. \url{https://doi.org/10.1111/j.1478-3231.2011.02501.x}.

\bibitem{ref034}
Jaeschke, H., and Ramachandran, A. (2024). Acetaminophen Hepatotoxicity: Paradigm for Understanding Mechanisms of Drug-Induced Liver Injury. Annu. Rev. Pathol. \textit{19}, 453--478. \url{https://doi.org/10.1146/annurev-pathmechdis-051122-094016}.

\bibitem{ref035}
Almadani, Y.H., Vorstenbosch, J., Davison, P.G., and Murphy, A.M. (2021). Wound Healing: A Comprehensive Review. Semin. Plast. Surg. \textit{35}, 141--144. \url{https://doi.org/10.1055/s-0041-1731791}.

\bibitem{ref036}
Bartold, M., and Ivanovski, S. (2025). Biological processes and factors involved in soft and hard tissue healing. Periodontol. 2000 \textit{97}, 16--42. \url{https://doi.org/10.1111/prd.12546}.

\bibitem{ref037}
Jin, C., Jin, Y., Ding, Z., Nuch, K.S., Han, M., Shim, J., Chien, P.N., and Heo, C.Y. (2025). Cellular and Molecular Mechanisms of Wound Repair: From Biology to Therapeutic Innovation. Cells \textit{14}, 1850. \url{https://doi.org/10.3390/cells14231850}.

\bibitem{ref038}
Bröhl, T., and Lehnertz, K. (2022). A straightforward edge centrality concept derived from generalizing degree and strength. Sci. Rep. \textit{12}, 4407. \url{https://doi.org/10.1038/s41598-022-08254-5}.

\bibitem{ref039}
Yuan, S., Liu, W., Yan, G., Tang, H., and Cheng, X. (2025). A Novel Method to Identify Important Edges in a Directed Network. Axioms \textit{14}, 258. \url{https://doi.org/10.3390/axioms14040258}.

\bibitem{ref040}
Setty, M., Kiseliovas, V., Levine, J., Gayoso, A., Mazutis, L., and Pe'er, D. (2019). Characterization of cell fate probabilities in single-cell data with Palantir. Nat. Biotechnol. \textit{37}, 451--460. \url{https://doi.org/10.1038/s41587-019-0068-4}.

\bibitem{ref041}
Lange, M., Bergen, V., Klein, M., Setty, M., Reuter, B., Bakhti, M., Lickert, H., Ansari, M., Schniering, J., Schiller, H.B., et al. (2022). CellRank for directed single-cell fate mapping. Nat. Methods \textit{19}, 159--170. \url{https://doi.org/10.1038/s41592-021-01346-6}.

\bibitem{ref042}
Breiman, L. (2001). Random Forests. Mach. Learn. \textit{45}, 5--32. \url{https://doi.org/10.1023/A:1010933404324}.

\bibitem{ref043}
Rossi, E., Charpentier, B., Di Giovanni, F., Frasca, F., Günnemann, S., and Bronstein, M.M. (2024). Edge Directionality Improves Learning on Heterophilic Graphs. In Proceedings of the Second Learning on Graphs Conference, Proceedings of Machine Learning Research 231, pp. 25:1--25:27. \url{https://proceedings.mlr.press/v231/rossi24a.html}.

\bibitem{ref044}
Kipf, T.N., and Welling, M. (2017). Semi-Supervised Classification with Graph Convolutional Networks. In International Conference on Learning Representations (ICLR). \url{https://doi.org/10.48550/arXiv.1609.02907}.

\bibitem{ref045}
Suarez-Lopez, L., Shui, B., Brubaker, D.K., Hill, M., Bergendorf, A., Changelian, P.S., Laguna, A., Starchenko, A., Lauffenburger, D.A., and Haigis, K.M. (2021). Cross-species transcriptomic signatures predict response to MK2 inhibition in mouse models of chronic inflammation. iScience \textit{24}, 103406. \url{https://doi.org/10.1016/j.isci.2021.103406}.

\bibitem{ref046}
Cao, G., Qiu, S., Hu, C., Kan, Y., Li, Z., Li, S., Feng, Q., Fu, W., and Wang, L. (2026). Cross-species scRNA-seq finds ideal mouse models for aortic dissection mechanisms. iScience \textit{29}, 115147. \url{https://doi.org/10.1016/j.isci.2026.115147}.

\bibitem{ref047}
Cao, J., Spielmann, M., Qiu, X., Huang, X., Ibrahim, D.M., Hill, A.J., Zhang, F., Mundlos, S., Christiansen, L., Steemers, F.J., et al. (2019). The single-cell transcriptional landscape of mammalian organogenesis. Nature \textit{566}, 496--502. \url{https://doi.org/10.1038/s41586-019-0969-x}.

\bibitem{ref048}
Chen, T., and Guestrin, C. (2016). XGBoost: A Scalable Tree Boosting System. In Proceedings of the 22nd ACM SIGKDD International Conference on Knowledge Discovery and Data Mining (ACM), pp. 785--794. \url{https://doi.org/10.1145/2939672.2939785}.

\bibitem{ref049}
Bhushan, B., and Apte, U. (2023). Regeneration and Recovery after Acetaminophen Hepatotoxicity. Livers \textit{3}, 300--309. \url{https://doi.org/10.3390/livers3020021}.

\bibitem{ref050}
Afrose, S., Song, W., Nemeroff, C.B., Lu, C., and Yao, D.D. (2022). Subpopulation-specific machine learning prognosis for underrepresented patients with double prioritized bias correction. Commun. Med. \textit{2}, 111. \url{https://doi.org/10.1038/s43856-022-00165-w}.

\bibitem{ref051}
Li, S., Zeng, W., Li, W., Liu, C.-C., Zhou, Y., Ni, X., Stackpole, M.L., Yeh, A.H., Melehy, A., Lu, D.S., et al. (2026). Reducing demographic bias in biomedical machine learning for cancer detection using cfDNA methylation. Genome Biol. \textit{27}, 141. \url{https://doi.org/10.1186/s13059-026-04006-0}.

\bibitem{ref052}
Wang, S.-W., Herriges, M.J., Hurley, K., Kotton, D.N., and Klein, A.M. (2022). CoSpar identifies early cell fate biases from single-cell transcriptomic and lineage information. Nat. Biotechnol. \textit{40}, 1066--1074. \url{https://doi.org/10.1038/s41587-022-01209-1}.

\bibitem{ref053}
Lange, M., Piran, Z., Klein, M., Spanjaard, B., Klein, D., Junker, J.P., Theis, F.J., and Nitzan, M. (2024). Mapping lineage-traced cells across time points with moslin. Genome Biol. \textit{25}, 277. \url{https://doi.org/10.1186/s13059-024-03422-4}.

\bibitem{ref054}
Fu, Y., Mathew, D., Wang, M., Chen, X.E., Lin, K.Z., Schaff, D., Shaffer, S.M., Pardoll, D.M., Jackson, C., and Zhang, N.R. (2025). Deciphering Cell Fate and Clonal Dynamics via Integrative Single-Cell Lineage Modeling. Preprint at Research Square. \url{https://doi.org/10.21203/rs.3.rs-7510946/v1}.

\bibitem{ref055}
Bonham-Carter, B., and Schiebinger, G. (2024). Cellular proliferation biases clonal lineage tracing and trajectory inference. Bioinformatics \textit{40}, btae483. \url{https://doi.org/10.1093/bioinformatics/btae483}.

\bibitem{ref056}
Aivazidis, A., Memi, F., Kleshchevnikov, V., Er, S., Clarke, B., Stegle, O., and Bayraktar, O.A. (2025). Cell2fate infers RNA velocity modules to improve cell fate prediction. Nat. Methods \textit{22}, 698--707. \url{https://doi.org/10.1038/s41592-025-02608-3}.

\bibitem{ref057}
Gu, Y., Song, Y., Blaauw, D., and Welch, J.D. (2026). Bayesian inference of RNA velocity incorporating timepoints, lineage bifurcations, and count data. PLoS Comput. Biol. \textit{22}, e1014060. \url{https://doi.org/10.1371/journal.pcbi.1014060}.

\bibitem{ref058}
Schleicher, J.T., Klingler, D., and Claassen, M. (2026). Accurate quantification of spliced and unspliced transcripts for single-cell RNA sequencing with tidesurf. PLoS One \textit{21}, e0355867. \url{https://doi.org/10.1371/journal.pone.0355867}.

\bibitem{ref059}
Chen, G., and Liu, Z.-P. (2022). Graph attention network for link prediction of gene regulations from single-cell RNA-sequencing data. Bioinformatics \textit{38}, 4522--4529. \url{https://doi.org/10.1093/bioinformatics/btac559}.

\bibitem{ref060}
Gan, Y., Yu, J., Xu, G., Yan, C., and Zou, G. (2024). Inferring gene regulatory networks from single-cell transcriptomics based on graph embedding. Bioinformatics \textit{40}, btae291. \url{https://doi.org/10.1093/bioinformatics/btae291}.

\bibitem{ref061}
Umbaugh, D.S., Nguyen, N.T., Smith, S.H., Ramachandran, A., and Jaeschke, H. (2024). The p21+ perinecrotic hepatocytes produce the chemokine CXCL14 after a severe acetaminophen overdose promoting hepatocyte injury and delaying regeneration. Toxicology \textit{504}, 153804. \url{https://doi.org/10.1016/j.tox.2024.153804}.

\bibitem{ref062}
Moris, N., Pina, C., and Martinez Arias, A. (2016). Transition states and cell fate decisions in epigenetic landscapes. Nat. Rev. Genet. \textit{17}, 693--703. \url{https://doi.org/10.1038/nrg.2016.98}.

\bibitem{ref063}
Zhang, J., Nie, Q., and Zhou, T. (2019). Revealing Dynamic Mechanisms of Cell Fate Decisions From Single-Cell Transcriptomic Data. Front. Genet. \textit{10}, 1280. \url{https://doi.org/10.3389/fgene.2019.01280}.

\bibitem{ref064}
Lotfollahi, M., Naghipourfar, M., Luecken, M.D., Khajavi, M., Büttner, M., Wagenstetter, M., Avsec, Ž., Gayoso, A., Yosef, N., Interlandi, M., et al. (2022). Mapping single-cell data to reference atlases by transfer learning. Nat. Biotechnol. \textit{40}, 121--130. \url{https://doi.org/10.1038/s41587-021-01001-7}.

\bibitem{ref065}
Michielsen, L., Lotfollahi, M., Strobl, D., Sikkema, L., Reinders, M.J.T., Theis, F.J., and Mahfouz, A. (2023). Single-cell reference mapping to construct and extend cell-type hierarchies. NAR Genom. Bioinform. \textit{5}, lqad070. \url{https://doi.org/10.1093/nargab/lqad070}.

\bibitem{ref066}
Rosen, Y., Brbić, M., Roohani, Y., Swanson, K., Li, Z., and Leskovec, J. (2024). Toward universal cell embeddings: integrating single-cell RNA-seq datasets across species with SATURN. Nat. Methods \textit{21}, 1492--1500. \url{https://doi.org/10.1038/s41592-024-02191-z}.

\bibitem{ref067}
Schächter, C., Hackenberg, M., Treppner, M., Raum, H., Bödecker, J., and Binder, H. (2025). scSpecies: enhancement of network architecture alignment in comparative single-cell studies. Genome Biol. \textit{26}, 397. \url{https://doi.org/10.1186/s13059-025-03866-2}.

\bibitem{ref068}
Chhatbar, C., Sankowski, R., Schulz, M., Shimizu, T., Schwabenland, M., Staszewski, O., Scheiwe, C., Nessler, S., Borst, K., Dumas, A.A., et al. (2026). A transcriptomic microglia taxonomy across mouse and human pathologies. Nat. Immunol. \textit{27}, 1066--1080. \url{https://doi.org/10.1038/s41590-026-02472-z}.

\bibitem{ref069}
Zhou, Y., Song, W.M., Andhey, P.S., Swain, A., Levy, T., Miller, K.R., Poliani, P.L., Cominelli, M., Grover, S., Gilfillan, S., et al. (2020). Human and mouse single-nucleus transcriptomics reveal TREM2-dependent and TREM2-independent cellular responses in Alzheimer's disease. Nat. Med. \textit{26}, 131--142. \url{https://doi.org/10.1038/s41591-019-0695-9}.

\bibitem{ref070}
Squair, J.W., Gautier, M., Kathe, C., Anderson, M.A., James, N.D., Hutson, T.H., Hudelle, R., Qaiser, T., Matson, K.J.E., Barraud, Q., et al. (2021). Confronting false discoveries in single-cell differential expression. Nat. Commun. \textit{12}, 5692. \url{https://doi.org/10.1038/s41467-021-25960-2}.

\bibitem{ref071}
Murphy, A.E., and Skene, N.G. (2022). A balanced measure shows superior performance of pseudobulk methods in single-cell RNA-sequencing analysis. Nat. Commun. \textit{13}, 7851. \url{https://doi.org/10.1038/s41467-022-35519-4}.

\bibitem{ref072}
Zimmerman, K.D., Espeland, M.A., and Langefeld, C.D. (2021). A practical solution to pseudoreplication bias in single-cell studies. Nat. Commun. \textit{12}, 738. \url{https://doi.org/10.1038/s41467-021-21038-1}.

\bibitem{ref073}
Kotulkar, M., Robarts, D.R., and Apte, U. (2023). HNF4\ensuremath{\alpha} in Hepatocyte Health and Disease. Semin. Liver Dis. \textit{43}, 234--244. \url{https://doi.org/10.1055/a-2097-0660}.

\bibitem{ref074}
Barbier-Torres, L., Chhimwal, J., Mato, J.M., and Lu, S.C. (2025). Methionine Adenosyltransferase 1A and S-Adenosylmethionine in Alcohol-Associated Liver Disease. Antioxidants \textit{14}, 1486. \url{https://doi.org/10.3390/antiox14121486}.

\bibitem{ref075}
Hyun, J., Oh, S.-H., Premont, R.T., Guy, C.D., Berg, C.L., and Diehl, A.M. (2019). Dysregulated activation of fetal liver programme in acute liver failure. Gut \textit{68}, 1076--1087. \url{https://doi.org/10.1136/gutjnl-2018-317603}.

\bibitem{ref076}
Ehle, C., Iyer-Bierhoff, A., Wu, Y., Xing, S., Kiehntopf, M., Mosig, A.S., Godmann, M., and Heinzel, T. (2024). Downregulation of HNF4A enables transcriptomic reprogramming during the hepatic acute-phase response. Commun. Biol. \textit{7}, 589. \url{https://doi.org/10.1038/s42003-024-06288-1}.

\bibitem{ref077}
Chen, G., Jia, H., Zhao, K., Gao, C., Liu, Y., Zhu, H., Yang, D., and Cai, Y. (2026). Spatial hepatology: Decoding liver zonation for metabolic and regenerative therapeutics. Int. J. Mol. Med. \textit{58}, 185. \url{https://doi.org/10.3892/ijmm.2026.5856}.

\bibitem{ref078}
Albadry, M., Küttner, J., Grzegorzewski, J., Dirsch, O., Kindler, E., Klopfleisch, R., Liska, V., Moulisova, V., Nickel, S., Palek, R., et al. (2024). Cross-species variability in lobular geometry and cytochrome P450 hepatic zonation: insights into CYP1A2, CYP2D6, CYP2E1 and CYP3A4. Front. Pharmacol. \textit{15}, 1404938. \url{https://doi.org/10.3389/fphar.2024.1404938}.

\bibitem{ref079}
Yasaka, T.M., Kim, C.K., Meadows, V., and Monga, S.P. (2026). Zonation, Zonation, Zonation: The Real Estate of the Liver. Annu. Rev. Pathol. \textit{21}, 185--212. \url{https://doi.org/10.1146/annurev-pathmechdis-042624-091820}.

\bibitem{ref080}
Wang, S., Wang, X., Shan, Y., Tan, Z., Su, Y., Cao, Y., Wang, S., Dong, J., Gu, J., and Wang, Y. (2024). Region-specific cellular and molecular basis of liver regeneration after acute pericentral injury. Cell Stem Cell \textit{31}, 341--358.e7. \url{https://doi.org/10.1016/j.stem.2024.01.013}.

\bibitem{ref081}
Wen, Y., and Ju, C. (2025). New insights into liver injury and regeneration from single-cell transcriptomics. eGastroenterology \textit{3}, e100202. \url{https://doi.org/10.1136/egastro-2025-100202}.

\bibitem{ref082}
Ruz Jurado, M., Tombor, L.S., Arsalan, M., Holubec, T., Emrich, F., Walther, T., Abplanalp, W., Fischer, A., Zeiher, A.M., Schulz, M.H., et al. (2024). Improved integration of single-cell transcriptome data demonstrates common and unique signatures of heart failure in mice and humans. GigaScience \textit{13}, giae011. \url{https://doi.org/10.1093/gigascience/giae011}.

\bibitem{ref083}
Seok, J., Warren, H.S., Cuenca, A.G., Mindrinos, M.N., Baker, H.V., Xu, W., Richards, D.R., McDonald-Smith, G.P., Gao, H., Hennessy, L., et al. (2013). Genomic responses in mouse models poorly mimic human inflammatory diseases. Proc. Natl. Acad. Sci. USA \textit{110}, 3507--3512. \url{https://doi.org/10.1073/pnas.1222878110}.

\bibitem{ref084}
Bishop, C.R., Dumenil, T., Rawle, D.J., Le, T.T., Yan, K., Tang, B., Hartel, G., and Suhrbier, A. (2022). Mouse models of COVID-19 recapitulate inflammatory pathways rather than gene expression. PLoS Pathog. \textit{18}, e1010867. \url{https://doi.org/10.1371/journal.ppat.1010867}.

\bibitem{ref085}
McGill, M.R., Williams, C.D., Xie, Y., Ramachandran, A., and Jaeschke, H. (2012). Acetaminophen-induced liver injury in rats and mice: comparison of protein adducts, mitochondrial dysfunction, and oxidative stress in the mechanism of toxicity. Toxicol. Appl. Pharmacol. \textit{264}, 387--394. \url{https://doi.org/10.1016/j.taap.2012.08.015}.

\bibitem{ref086}
Williams, C.D., Bajt, M.L., Sharpe, M.R., McGill, M.R., Farhood, A., and Jaeschke, H. (2014). Neutrophil activation during acetaminophen hepatotoxicity and repair in mice and humans. Toxicol. Appl. Pharmacol. \textit{275}, 122--133. \url{https://doi.org/10.1016/j.taap.2014.01.004}.

\bibitem{ref087}
Dixit, A., Parnas, O., Li, B., Chen, J., Fulco, C.P., Jerby-Arnon, L., Marjanovic, N.D., Dionne, D., Burks, T., Raychowdhury, R., et al. (2016). Perturb-Seq: Dissecting Molecular Circuits with Scalable Single-Cell RNA Profiling of Pooled Genetic Screens. Cell \textit{167}, 1853--1866.e17. \url{https://doi.org/10.1016/j.cell.2016.11.038}.

\bibitem{ref088}
Bouman, B.J., Demerdash, Y., Sood, S., Grünschläger, F., Pilz, F., Itani, A.R., Kuck, A., Marot-Lassauzaie, V., Haas, S., Haghverdi, L., et al. (2024). Single-cell time series analysis reveals the dynamics of HSPC response to inflammation. Life Sci. Alliance \textit{7}, e202302309. \url{https://doi.org/10.26508/lsa.202302309}.

\bibitem{ref089}
Wahle, P., Brancati, G., Harmel, C., He, Z., Gut, G., del Castillo, J.S., Xavier da Silveira dos Santos, A., Yu, Q., Noser, P., Fleck, J.S., et al. (2023). Multimodal spatiotemporal phenotyping of human retinal organoid development. Nat. Biotechnol. \textit{41}, 1765--1775. \url{https://doi.org/10.1038/s41587-023-01747-2}.

\bibitem{ref090}
Peidli, S., Green, T.D., Shen, C., Gross, T., Min, J., Garda, S., Yuan, B., Schumacher, L.J., Taylor-King, J.P., Marks, D.S., et al. (2024). scPerturb: harmonized single-cell perturbation data. Nat. Methods \textit{21}, 531--540. \url{https://doi.org/10.1038/s41592-023-02144-y}.

\bibitem{ref091}
Ma, A., Wang, X., Li, J., Wang, C., Xiao, T., Liu, Y., Cheng, H., Wang, J., Li, Y., Chang, Y., et al. (2023). Single-cell biological network inference using a heterogeneous graph transformer. Nat. Commun. \textit{14}, 964. \url{https://doi.org/10.1038/s41467-023-36559-0}.

\bibitem{ref092}
Wang, J., Chen, Y., and Zou, Q. (2023). Inferring gene regulatory network from single-cell transcriptomes with graph autoencoder model. PLoS Genet. \textit{19}, e1010942. \url{https://doi.org/10.1371/journal.pgen.1010942}.

\bibitem{ref093}
Kommu, S., Wang, Y., Wang, Y., and Wang, X. (2025). Prediction of gene regulatory connections with joint single-cell foundation models and graph-based learning. Bioinformatics \textit{41}, i619--i627. \url{https://doi.org/10.1093/bioinformatics/btaf217}.

\bibitem{ref094}
Li, G., Qiao, R., Xu, Y., Jiao, F., Li, P., and Chen, L. (2026). Inferring cell-specific gene regulatory networks based on causal graph embedding. Cell Rep. Methods \textit{6}, 101423. \url{https://doi.org/10.1016/j.crmeth.2026.101423}.

\bibitem{ref095}
Malagoli, G., Hanel, P., Danese, A., Wolf, G., and Colomé-Tatché, M. (2026). Geometry-aware graph attention networks to explain single-cell chromatin states and gene expression with SEAGALL. Genome Biol. \textit{27}, 188. \url{https://doi.org/10.1186/s13059-026-04066-2}.

\bibitem{ref096}
Wang, Y., Tian, S., and Li, D. (2026). IAGRN: An Interleaved-Attention Graph Neural Network for Gene Regulatory Network Inference. Int. J. Mol. Sci. \textit{27}, 7300. \url{https://doi.org/10.3390/ijms27167300}.

\bibitem{ref097}
Dip, S.A., and Zhang, L. (2026). Predicting Unseen Gene Perturbation Response Using Graph Neural Networks with Biological Priors. Preprint at bioRxiv. \url{https://doi.org/10.64898/2026.03.23.713780}.

\bibitem{ref098}
Stuart, T., Butler, A., Hoffman, P., Hafemeister, C., Papalexi, E., Mauck, W.M., III, Hao, Y., Stoeckius, M., Smibert, P., and Satija, R. (2019). Comprehensive Integration of Single-Cell Data. Cell \textit{177}, 1888--1902.e21. \url{https://doi.org/10.1016/j.cell.2019.05.031}.

\bibitem{ref099}
Kang, J.B., Nathan, A., Weinand, K., Zhang, F., Millard, N., Rumker, L., Moody, D.B., Korsunsky, I., and Raychaudhuri, S. (2021). Efficient and precise single-cell reference atlas mapping with Symphony. Nat. Commun. \textit{12}, 5890. \url{https://doi.org/10.1038/s41467-021-25957-x}.

\bibitem{ref100}
Wu, H., Kirita, Y., Donnelly, E.L., and Humphreys, B.D. (2019). Advantages of Single-Nucleus over Single-Cell RNA Sequencing of Adult Kidney: Rare Cell Types and Novel Cell States Revealed in Fibrosis. J. Am. Soc. Nephrol. \textit{30}, 23--32. \url{https://doi.org/10.1681/ASN.2018090912}.

\bibitem{ref101}
Jaeschke, H., Xie, Y., and McGill, M.R. (2014). Acetaminophen-induced Liver Injury: from Animal Models to Humans. J. Clin. Transl. Hepatol. \textit{2}, 153--161. \url{https://doi.org/10.14218/JCTH.2014.00014}.

\bibitem{ref102}
Yang, X., Greenhaw, J., Ali, A., Shi, Q., Roberts, D.W., Hinson, J.A., Muskhelishvili, L., Beger, R., Pence, L.M., Ando, Y., et al. (2012). Changes in mouse liver protein glutathionylation after acetaminophen exposure. J. Pharmacol. Exp. Ther. \textit{340}, 360--368. \url{https://doi.org/10.1124/jpet.111.187948}.

\bibitem{ref103}
Wang, Q., Li, W.-X., Dai, S.-X., Guo, Y.-C., Han, F.-F., Zheng, J.-J., Li, G.-H., and Huang, J.-F. (2017). Meta-Analysis of Parkinson's Disease and Alzheimer's Disease Revealed Commonly Impaired Pathways and Dysregulation of NRF2-Dependent Genes. J. Alzheimers Dis. \textit{56}, 1525--1539. \url{https://doi.org/10.3233/JAD-161032}.

\bibitem{ref104}
Kanehisa, M., Furumichi, M., Sato, Y., Ishiguro-Watanabe, M., and Tanabe, M. (2021). KEGG: integrating viruses and cellular organisms. Nucleic Acids Res. \textit{49}, D545--D551. \url{https://doi.org/10.1093/nar/gkaa970}.

\bibitem{ref105}
Ritchie, S.C., Watts, S., Fearnley, L.G., Abraham, G., and Inouye, M. (2016). A Scalable Permutation Approach Reveals Replication and Preservation Patterns of Network Modules in Large Datasets. Cell Syst. \textit{3}, 71--82. \url{https://doi.org/10.1016/j.cels.2016.06.012}.

\bibitem{ref106}
Fulcher, B.D., Arnatkevičiūtė, A., and Fornito, A. (2021). Overcoming false-positive gene-category enrichment in the analysis of spatially resolved transcriptomic brain atlas data. Nat. Commun. \textit{12}, 2669. \url{https://doi.org/10.1038/s41467-021-22862-1}.

\bibitem{ref107}
Kana, O., Nault, R., Filipovic, D., Marri, D., Zacharewski, T., and Bhattacharya, S. (2023). Generative modeling of single-cell gene expression for dose-dependent chemical perturbations. Patterns \textit{4}, 100817. \url{https://doi.org/10.1016/j.patter.2023.100817}.

\bibitem{ref108}
Wang, G., Liu, T., Zhao, J., Cheng, Y., and Zhao, H. (2024). Modeling and predicting single-cell multi-gene perturbation responses with scLAMBDA. Preprint at bioRxiv. \url{https://doi.org/10.1101/2024.12.04.626878}.

\bibitem{ref109}
Wolf, F.A., Angerer, P., and Theis, F.J. (2018). SCANPY: large-scale single-cell gene expression data analysis. Genome Biol. \textit{19}, 15. \url{https://doi.org/10.1186/s13059-017-1382-0}.

\bibitem{ref110}
McInnes, L., Healy, J., Saul, N., and Großberger, L. (2018). UMAP: Uniform Manifold Approximation and Projection. J. Open Source Softw. \textit{3}, 861. \url{https://doi.org/10.21105/joss.00861}.

\bibitem{ref111}
Traag, V.A., Waltman, L., and van Eck, N.J. (2019). From Louvain to Leiden: guaranteeing well-connected communities. Sci. Rep. \textit{9}, 5233. \url{https://doi.org/10.1038/s41598-019-41695-z}.

\bibitem{ref112}
Lopez, R., Regier, J., Cole, M.B., Jordan, M.I., and Yosef, N. (2018). Deep generative modeling for single-cell transcriptomics. Nat. Methods \textit{15}, 1053--1058. \url{https://doi.org/10.1038/s41592-018-0229-2}.

\bibitem{ref113}
Gayoso, A., Lopez, R., Xing, G., Boyeau, P., Valiollah Pour Amiri, V., Hong, J., Wu, K., Jayasuriya, M., Mehlman, E., Langevin, M., et al. (2022). A Python library for probabilistic analysis of single-cell omics data. Nat. Biotechnol. \textit{40}, 163--166. \url{https://doi.org/10.1038/s41587-021-01206-w}.

\bibitem{ref114}
Trapnell, C., Cacchiarelli, D., Grimsby, J., Pokharel, P., Li, S., Morse, M., Lennon, N.J., Livak, K.J., Mikkelsen, T.S., and Rinn, J.L. (2014). The dynamics and regulators of cell fate decisions are revealed by pseudotemporal ordering of single cells. Nat. Biotechnol. \textit{32}, 381--386. \url{https://doi.org/10.1038/nbt.2859}.

\bibitem{ref115}
Qiu, X., Mao, Q., Tang, Y., Wang, L., Chawla, R., Pliner, H.A., and Trapnell, C. (2017). Reversed graph embedding resolves complex single-cell trajectories. Nat. Methods \textit{14}, 979--982. \url{https://doi.org/10.1038/nmeth.4402}.

\bibitem{ref116}
McInnes, L., Healy, J., and Melville, J. (2018). UMAP: Uniform Manifold Approximation and Projection for Dimension Reduction. Preprint at arXiv (version 3, revised 2020; first posted 2018). \url{https://doi.org/10.48550/arXiv.1802.03426}.

\bibitem{ref117}
Pedregosa, F., Varoquaux, G., Gramfort, A., Michel, V., Thirion, B., Grisel, O., Blondel, M., Prettenhofer, P., Weiss, R., Dubourg, V., et al. (2011). Scikit-learn: Machine Learning in Python. J. Mach. Learn. Res. \textit{12}, 2825--2830. \url{https://www.jmlr.org/papers/v12/pedregosa11a.html}.

\bibitem{ref118}
Fey, M., and Lenssen, J.E. (2019). Fast Graph Representation Learning with PyTorch Geometric. In ICLR Workshop on Representation Learning on Graphs and Manifolds. \url{https://doi.org/10.48550/arXiv.1903.02428}.

\bibitem{ref119}
Paszke, A., Gross, S., Massa, F., Lerer, A., Bradbury, J., Chanan, G., Killeen, T., Lin, Z., Gimelshein, N., Antiga, L., et al. (2019). PyTorch: An Imperative Style, High-Performance Deep Learning Library. In Advances in Neural Information Processing Systems 32, H. Wallach, H. Larochelle, A. Beygelzimer, F. d'Alché-Buc, E. Fox, and R. Garnett, eds. (Curran Associates, Inc.), pp. 8024--8035. \url{https://papers.neurips.cc/paper_files/paper/2019/hash/bdbca288fee7f92f2bfa9f7012727740-Abstract.html}.

\bibitem{ref120}
Kingma, D.P., and Ba, J. (2015). Adam: A Method for Stochastic Optimization. In International Conference on Learning Representations (ICLR). \url{https://doi.org/10.48550/arXiv.1412.6980}.

\bibitem{ref121}
Smedley, D., Haider, S., Durinck, S., Pandini, L., Provero, P., Allen, J., Arnaiz, O., Awedh, M.H., Baldock, R., Barbiera, G., et al. (2015). The BioMart community portal: an innovative alternative to large, centralized data repositories. Nucleic Acids Res. \textit{43}, W589--W598. \url{https://doi.org/10.1093/nar/gkv350}.

\bibitem{ref122}
Virtanen, P., Gommers, R., Oliphant, T.E., Haberland, M., Reddy, T., Cournapeau, D., Burovski, E., Peterson, P., Weckesser, W., Bright, J., et al. (2020). SciPy 1.0: fundamental algorithms for scientific computing in Python. Nat. Methods \textit{17}, 261--272. \url{https://doi.org/10.1038/s41592-019-0686-2}.

\bibitem{ref123}
Kuleshov, M.V., Jones, M.R., Rouillard, A.D., Fernandez, N.F., Duan, Q., Wang, Z., Koplev, S., Jenkins, S.L., Jagodnik, K.M., Lachmann, A., et al. (2016). Enrichr: a comprehensive gene set enrichment analysis web server 2016 update. Nucleic Acids Res. \textit{44}, W90--W97. \url{https://doi.org/10.1093/nar/gkw377}.

\bibitem{ref124}
Fang, Z., Liu, X., and Peltz, G. (2023). GSEApy: a comprehensive package for performing gene set enrichment analysis in Python. Bioinformatics \textit{39}, btac757. \url{https://doi.org/10.1093/bioinformatics/btac757}.

\bibitem{ref125}
Ashburner, M., Ball, C.A., Blake, J.A., Botstein, D., Butler, H., Cherry, J.M., Davis, A.P., Dolinski, K., Dwight, S.S., Eppig, J.T., et al. (2000). Gene Ontology: tool for the unification of biology. Nat. Genet. \textit{25}, 25--29. \url{https://doi.org/10.1038/75556}.

\bibitem{ref126}
Imani, M., Joudaki, M., Bagheri, A., and Arabnia, H.R. (2026). Why ROC-AUC Is Misleading for Highly Imbalanced Data: In-Depth Evaluation of MCC, F2-Score, H-Measure, and AUC-Based Metrics Across Diverse Classifiers. Technologies \textit{14}, 54. \url{https://doi.org/10.3390/technologies14010054}.

\end{thebibliography}
\endgroup
